\subsection{Quotient spaces and subspaces of a weak space}

Let F, G be two real vector spaces in duality. Let M be a vector subspace of F, and consider the subspace N of the orthogonal \(M^{\circ}\) in G; if \(y_{1}, y_{2}\) are two points of G that are congruent mod. N then \(\langle x, y_{1} \rangle = \langle x, y_{2} \rangle\) for all \(x \in M\) . For each class \(\dot{y}\) mod. N, denote the common value of \(\langle x, y \rangle\) when y varies in \(\dot{y}\) by \(\langle x, \dot{y} \rangle\) ; clearly \((x, \dot{y}) \mapsto \langle x, \dot{y} \rangle\) is a bilinear form on \(M \times (G/N)\) .

PROPOSITION 7. --- Let M be a vector subspace of F and N a vector subspace of G where F and G are two vector spaces in duality. Suppose that M and N are orthogonal (which is equivalent to saying that \(N \subset M^{\circ}\) , or \(M \subset N^{\circ}\) ). The vector spaces M and G/N are then in duality by the bilinear form \((x, \dot{y}) \mapsto \langle x, \dot{y} \rangle\) .

\begin{enumerate}
\def\labelenumi{(\roman{enumi})}
\item
  The topology \(\sigma(\mathbf{M},\mathbf{G} / \mathbf{N})\) for this duality is induced by \(\sigma(\mathbf{F},\mathbf{G})\) (and in particular we have \(\sigma(\mathbf{F},\mathbf{G}) = \sigma(\mathbf{F},\mathbf{G} / \mathbf{F}^{\circ})\) ).
\item
  The topology \(\sigma(\mathbf{G}/\mathbf{N}, \mathbf{M})\) for this duality is coarser than the quotient topology of \(\sigma(\mathbf{G}, \mathbf{F})\) by \(\mathbf{N}\) ; these topologies are identical if and only if \(\mathbf{M} + \mathbf{G}^{\circ} = \mathbf{N}^{\circ}\) .
\setcounter{enumi}{0}
\item
  Every element of \(\mathbf{G} / \mathbf{N}\) is a class mod. \(\mathbf{N}\) of an element of \(\mathbf{G}\) ; if \(z_{i}\) ( \(1 \leqslant i \leqslant n\) ) are elements of \(\mathbf{G}\) and \(\dot{z}_{i}\) ( \(1 \leqslant i \leqslant n\) ) is the class of \(z_{i}\) in \(\mathbf{G} / \mathbf{N}\) then the set of \(y \in \mathbf{M}\) such that \(|\langle y, \dot{z}_{i} \rangle| \leqslant \alpha\) for \(1 \leqslant i \leqslant n\) is the trace on \(\mathbf{M}\) of the set of those \(x \in \mathbf{F}\) such that \(|\langle x, z_{i} \rangle| \leqslant \alpha\) for \(1 \leqslant i \leqslant n\) ; the conclusion follows from the definition of neighbourhoods of 0 for the weak topology.
\item
  Let \(p: \mathbf{G} \to \mathbf{G} / \mathbf{N}\) be the canonical surjection. We show that the quotient topology \(\mathcal{T}\) of \(\sigma(\mathbf{G}, \mathbf{F})\) by \(\mathbf{N}\) is identical with \(\sigma(\mathbf{G} / \mathbf{N}, \mathbf{N}^{\circ})\) . As, for \(z \in \mathbf{G}\) , \(y \in \mathbf{N}^{\circ}\) , we have \(\langle y, p(z) \rangle = \langle y, z \rangle\) , it follows that every neighbourhood of 0 for \(\sigma(\mathbf{G} / \mathbf{N}, \mathbf{N}^{\circ})\) is of the form \(p(\mathbf{V})\) , where \(\mathbf{V}\) is a neighbourhood of 0 for \(\sigma(\mathbf{G}, \mathbf{F})\) saturated for the relation \(z - z' \in \mathbf{N}\) , therefore \(\mathcal{T}\) is finer than \(\sigma(\mathbf{G} / \mathbf{N}, \mathbf{N}^{\circ})\) . Conversely let \(\mathbf{U} = \mathbf{W}(y_1, ..., y_n; \alpha)\) be a neighbourhood of 0 in \(\mathbf{G}\) for \(\sigma(\mathbf{G}, \mathbf{F})\) , where \(y_i \in \mathbf{F}\) for \(1 \leqslant i \leqslant n\) and \(\alpha > 0\) ; we are going to see that for \(1 \leqslant i \leqslant n\) , there exist elements \(t_i \in \mathbf{N}^{\circ}\) such that if one puts \(\mathbf{U}' = \mathbf{W}(t_1, ..., t_n; \alpha)\) , then \(p(\mathbf{U}') \subset p(\mathbf{U})\) ; this will show that \(\sigma(\mathbf{G} / \mathbf{N}, \mathbf{N}^{\circ})\) is finer than \(\mathcal{T}\) and therefore is actually identical with \(\mathcal{T}\) . Now, let \(L\) be the vector subspace of \(F\) generated by \(N^{\circ}\) and the \(y_i\) , and denote by \(P\) the complementary subspace of \(N^{\circ}\) in \(L\) ; it is of finite dimension, say \(m\) . Let \((x_j)_{1 \leqslant j \leqslant m}\) be a basis of \(P\) ; the restrictions to \(N\) of the linear forms \(x \mapsto \langle x_j, z \rangle\) are linearly independent, since otherwise there exists \(x \neq 0\) in \(P\) such that \(\langle x, z \rangle = 0\) for all \(z \in N\) , that is to say \(x \in N^{\circ}\) , which contradicts the definition of \(P\) . Thus we conclude that for all \(z' \in G\) , there exists \(s \in N\) such that \(\langle x_j, z' \rangle = \langle x_j, s \rangle\) for all \(j\) ; if \(z' = z + s\) , we have \(\langle x, z \rangle = 0\) for all \(x \in P\) . This being so, put \(y_i = t_i + w_i\) , where \(t_i \in N^{\circ}\) and \(w_i \in P\) ; we have \(\langle y_i, z \rangle = \langle t_i, z \rangle = \langle t_i, z' \rangle\) for \(1 \leqslant i \leqslant n\) ; therefore, for all \(z' \in U'\) , there exists \(z \in U\) such that \(z' - z \in N\) , that is to say we have \(p(\mathrm{U}') \subset p(\mathrm{U})\) .
\end{enumerate}

Returning to the case where M is any subspace of \(N^{\circ}\) , note that evidently \(\sigma(\mathrm{G}/\mathrm{N},\mathrm{M})=\sigma(\mathrm{G}/\mathrm{N},\mathrm{M}+\mathrm{G}^{\circ})\) ; further, from prop. 3 of II, p.~43, we see that, if \(y\in N^{\circ}\) is such that the linear form \(\dot{z}\mapsto\langle y,\dot{z}\rangle\) is continuous for \(\sigma(\mathrm{G}/\mathrm{N},\mathrm{M})\) , then necessarily \(y\in M+G^{\circ}\) . We conclude that the condition \(M+G^{\circ}=N^{\circ}\) is necessary and sufficient for the quotient topology T to be equal to \(\sigma(\mathrm{G}/\mathrm{N},\mathrm{M})\) .

Remark. --- The duality between M and G/N (where M and N are two orthogonal subspaces) is separating in M, if and only if \(M \cap G^{\circ} = \{0\}\) ; it is separating in G/N, if and only if \(N = M^{\circ}\) .

COROLLARY 1. --- Suppose that the duality between F and G is separating in F. For a vector subspace M of F the topology \(\sigma(\mathbf{G}/\mathbf{M}^{\circ}, \mathbf{M})\) is identical with quotient topology of \(\sigma(\mathbf{G}, \mathbf{F})\) by \(\mathbf{M}^{\circ}\) , if and only if M is closed for the topology \(\sigma(\mathbf{F}, \mathbf{G})\) .

This follows from prop. 7 putting \(\mathbf{N} = \mathbf{M}^{\circ}\) , and recalling that \(\mathbf{M}^{\circ \circ}\) is the closure of \(\mathbf{M}\) for \(\sigma (\mathrm{F},\mathrm{G})\) (II, p.~45, cor. 2).

COROLLARY 2. --- If M is of finite dimension n and the duality is separating in F, then \(M^{\circ}\) is of codimension n in G. If M is closed for \(\sigma(F, G)\) and of finite codimension n and if the duality is separating in G, then \(M^{\circ}\) is of dimension n.

For, \(G/M^{\circ}\) is in separating duality with M; if M is of dimension n, the same is therefore true of \(G/M^{\circ}\) (II, p.~41, example 1). If M is closed, \(F/M = F/M^{\circ\circ}\) is in separating duality with \(M^{\circ}\) ; if F/M is of dimension n, it is therefore the same for \(M^{\circ}\) (II, p.~41, example 1).

COROLLARY 3. --- Let (F, G), (F₁, G₁) be two pairs of spaces in separating duality and let u be a linear mapping of F in F₁, which is continuous for σ(F, G) and σ(F₁, G₁). Then u is a strict morphism of F in F₁, if and only if, Im(\({}^t u\)) is a closed subspace in G for σ(G, F).

Let \(\mathbf{N} = \operatorname{Im}(^t u) \subset \mathbf{G}\) ; we know that \(\mathrm{N}^{\circ} = \mathrm{Ker}(u)\) in \(\mathbf{F}\) (II, p.~47, formula (3)). Let \(p: \mathbf{F} \to \mathbf{F}/\mathbf{N}^{\circ}\) be the canonical mapping so that \(u\) factorises as

\[
u: \mathrm{F} \xrightarrow {p} \mathrm{F} / \mathrm{N} ^ {\circ} \xrightarrow {w} \mathrm{F} _ {1},
\]

where w is injective. The spaces \(F/N^{\circ}\) and N are in separating duality and by formula (1) of II, p.~48, we have \(\langle w(\dot{y}), z_{1} \rangle = \langle \dot{y}, {}^{t}u(z_{1}) \rangle\) for all \(\dot{y} \in F/N^{\circ}\) and \(z_{1} \in G_{1}\) . This relation shows that w is an isomorphism of \(F/N^{\circ}\) , carrying the topology \(\sigma(F/N^{\circ}, N)\) , on \(u(F)\) with the topology induced by \(\sigma(F_{1}, G_{1})\) . The conclusion results therefore from cor. 1 and the definition of a strict morphism.

COROLLARY 4. --- Let \((\mathrm{F}, \mathrm{G})\) , \((\mathrm{F}_{1}, \mathrm{G}_{1})\) be two pairs in separating duality, and let u be a linear mapping of F in \(F_{1}\) that is continuous for \(\sigma(\mathrm{F}, \mathrm{G})\) and \(\sigma(\mathrm{F}_{1}, \mathrm{G}_{1})\) . Then u is surjective, if and only if, \(^{t}u\) is an isomorphism of \(G_{1}\) (with topology \(\sigma(G_{1}, F_{1})\) ) on \(^{t}u(G_{1})\) with the topology induced by \(\sigma(\mathrm{G}, \mathrm{F})\) .

For, to say that \(u(\mathrm{F}) = \mathrm{F}_{1}\) is equivalent to saying that \(u(\mathrm{F})\) is closed and everywhere dense in \(F_{1}\) for \(\sigma(\mathrm{F}_{1}, \mathrm{G}_{1})\) ; cor. 4 follows then from cor. 3 applied to \(^{t}u\) and of II, p.~47, cor. 2.

Remarks. --- 1) Let \((\mathrm{F}_1, \mathrm{G}_1), (\mathrm{F}_2, \mathrm{G}_2), (\mathrm{F}_3, \mathrm{G}_3)\) be three pairs of spaces in separating duality and consider a sequence of two linear mappings

\[
\mathrm{F} _ {1} \xrightarrow {u} \mathrm{F} _ {2} \xrightarrow {v} \mathrm{F} _ {3}\tag{5}
\]

that are continuous for the weak topologies corresponding respectively with \(\mathbf{G}_1\) , \(\mathbf{G}_2\) , \(\mathbf{G}_3\) ; we consider the sequence of transposed mappings

\[
\mathrm{G} _ {3} \stackrel {{t _ {v}}} {{\to}} \mathrm{G} _ {2} \stackrel {{t _ {u}}} {{\to}} \mathrm{G} _ {1}.\tag{6}
\]

It is clear that \({}^t (v\circ u) = {}^t u\circ {}^tv\) , therefore the relation \(v\circ u = 0\) is equivalent to \({}^t u\circ {}^tv = 0\) . The sequence (5) is exact if, and only if, the three following conditions are satisfied a) \({}^t u\circ {}^tv = 0\) :

\begin{enumerate}
\def\labelenumi{\alph{enumi})}
\setcounter{enumi}{1}
\item
  \(\operatorname{Im}(^{t}v)\) is dense in \(\operatorname{Ker}(^{t}u)\) ;
\item
  \(^t u\) is a strict morphism of \(G_2\) in \(G_1\) .
\end{enumerate}

This follows in effect from cor. 3 of II, p.~49 and formulae (3) and (4) of II, p.~47. 2) It must not be thought that when \(u\) is a strict morphism of \(F\) in \(F_1\) , then \(^t u\) is necessarily a strict morphism of \(G_1\) in \(G\) ; in other words \(u\) can be a strict morphism without \(u(F)\) being closed in \(F_1\) for \(\sigma(F_1, G_1)\) . This is shown by the example where \(F\) is a non-closed subspace of \(F_1\) and \(G = G_1 / F^\circ\) , \(u\) being the canonical injection. Similarly, the fact that the sequence (5) is exact does not necessarily imply that (6) is exact, however, if the sequence (5) is exact and if \(v\) is a strict morphism, then the sequence (6) is exact, by the remark 1 and by II, p.~49, cor. 3.

\subsection{Products of weak topologies}

PROPOSITION 8. --- Let \((\mathbf{F}_{\iota}, \mathbf{G}_{\iota})_{\iota \in \mathbf{I}}\) be a family of pairs of spaces in duality. Let \(\mathbf{F} = \prod_{\iota \in \mathbf{I}} \mathbf{F}_{\iota}\) be the product space of the \(\mathbf{F}_{\iota}\) and \(\mathbf{G} = \bigoplus_{\iota \in \mathbf{I}} \mathbf{G}_{\iota}\) be the direct sum of the \(\mathbf{G}_{\iota}\) . If, for all \(x = (x_{\iota}) \in \mathbf{F}\) and all \(y = (y_{\iota}) \in \mathbf{G}\) , we write \(\langle x, y \rangle = \sum_{\iota \in \mathbf{I}} \langle x_{\iota}, y_{\iota} \rangle\) (a sum which has only finitely many non-zero terms) then the topology \(\sigma(\mathbf{F}, \mathbf{G})\) (relative to the bilinear form \((x, y) \mapsto \langle x, y \rangle\) ) is the product of the topologies \(\sigma(\mathbf{F}_{\iota}, \mathbf{G}_{\iota})\) .

For, given a topology \(\mathcal{T}\) on \(\mathbf{F}\) ; in order that, for all \(y \in \mathbf{G}\) , the linear form \(x \mapsto \langle x, y \rangle\) should be continuous for \(\mathcal{T}\) , it is necessary and sufficient, by the definition of \(\langle x, y \rangle\) , that each of the mappings \(x \mapsto \langle \mathrm{pr}_i x, y_i \rangle\) should be continuous for \(\mathcal{T}\) , where \(i\) is arbitrary in I and \(y_i\) in \(\mathbf{G}_i\) ; but this means that each of the mappings \(\mathrm{pr}_i\) of F in \(\mathbf{F}_i\) is continuous for \(\mathcal{T}\) and for \(\sigma(\mathbf{F}_i, \mathbf{G}_i)\) (I, p.~10, cor. 1); this completes the demonstration.

Remark. --- The duality between F and G is separating in F (resp. in G) if and only if for all \(\iota\in I\) , the duality between \(F_{\iota}\) and \(G_{\iota}\) is separating in \(F_{\iota}\) (resp. in \(G_{\iota}\) ). If the duality between F and G is separating in F (resp. G), then, in F (resp. G), the subspace orthogonal to one \(G_{\iota}\) (resp. \(F_{\iota}\) ), canonically identified with a subspace of G (resp. F) is the subspace of the product of the \(F_{\kappa}\) where \(\kappa\neq\iota\) (resp. the direct sum of the \(G_{\kappa}\) such that \(\kappa\neq\iota\) ).

COROLLARY 1. --- Let F and G be two vector spaces in separating duality. If the space F (with \(\sigma(F, G)\) ) is the direct topological sum of two subspaces M, N then the space G (with \(\sigma(G, F)\) ) is the direct topological sum of the subspaces M°, N° orthogonal respectively to M and N.

For let \(p: \mathbf{F} \to \mathbf{M}\) , \(q: \mathbf{F} \to \mathbf{N}\) be the projectors corresponding to the decomposition of \(\mathbf{F}\) into the direct sum of \(\mathbf{M}\) and \(\mathbf{N}\) ; in these conditions the mapping \((p, q): \mathbf{F} \to \mathbf{M} \times \mathbf{N}\) is a topological isomorphism. If \(\mathbf{M}_1 = \mathbf{G} / \mathbf{M}^\circ\) , \(\mathbf{N}_1 = \mathbf{G} / \mathbf{N}^\circ\) , then the topologies on \(\mathbf{M}\) and \(\mathbf{N}\) (induced by that of \(\mathbf{F}\) ) are identical with \(\sigma(\mathbf{M}, \mathbf{M}_1)\) , \(\sigma(\mathbf{N}, \mathbf{N}_1)\) respectively (II, p.~48, prop. 7). The mapping \(^t(p, q): \mathbf{M}_1 \times \mathbf{N}_1 \to \mathbf{G}\) is a topological isomorphism when we give \(\mathbf{M}_1, \mathbf{N}_1\) and \(\mathbf{G}\) the topologies \(\sigma(\mathbf{M}_1, \mathbf{M})\) , \(\sigma(\mathbf{N}_1, \mathbf{N})\) and \(\sigma(\mathbf{G}, \mathbf{F})\) , by prop. 8. Under this mapping \(\mathbf{M}_1\) (resp. \(\mathbf{N}_1\) ) has as its image in \(\mathbf{G}\) the subspace \(\mathbf{N}^\circ\) (resp. \(\mathbf{M}^\circ\) ), and the topology \(\sigma(\mathbf{M}_1, \mathbf{M})\) (resp. \(\sigma(\mathbf{N}_1, \mathbf{N})\) ) has as its image the topology induced on \(\mathbf{N}^\circ\) (resp. \(\mathbf{M}^\circ\) ) by \(\sigma(\mathbf{G}, \mathbf{F})\) , from which the corollary follows.

COROLLARY 2. --- Let \((e_{\iota})_{\iota \in \mathbb{J}}\) be a basis of the vector space F with dual F*, and let \(u: \mathbf{R}^{(I)} \to \mathbf{F}\) be an (algebraic) isomorphism defined by this basis. Then the transposed mapping \({}^t u: \mathbf{F}^* \to \mathbf{R}^I\) is a topological isomorphism when F* carries the topology \(\sigma(\mathbf{F}^*, \mathbf{F})\) and \(\mathbf{R}^I\) the product topology.

We know (A, II, § 2.6, prop. 10) that \({}^t u\) is a bijection, and that if for a \(x^* \in \mathbf{F}^*\) , we put \(\langle e_i, x^* \rangle = \xi_i^*\) for all \(i \in I\) , then the image \({}^t u(x^*)\) is the vector \((\xi_i^*)\) of \(\mathbf{R}^I\) , so that, for all \(x = \sum_{i} \xi_i e_i\) in \(F\) , we have \(\langle x, x^* \rangle = \sum_{i \in I} \xi_i \xi_i^*\) . The corollary then follows from this formula and prop. 8.

\subsection{Weakly complete spaces}

PROPOSITION 9. --- Let F, G be two vector spaces in separating duality. If \(\hat{\mathbf{F}}\) is the completion of the space F for the topology \(\sigma(\mathrm{F},\mathrm{G})\) and if we consider the canonical injection \(j:\mathrm{F}\to\mathrm{G}^{*}\) , where \(\mathrm{G}^{*}\) has the topology \(\sigma(\mathrm{G}^{*},\mathrm{G})\) , then the continuous extension \(\hat{j}:\hat{\mathrm{F}}\to\mathrm{G}^{*}\) of \(j\) is an isomorphism of topological vector spaces.

For, we see that \(G^{*}\) , endowed with \(\sigma(G^{*}, G)\) , is Hausdorff and complete (II, p.~51, cor. 2); if we identify F by j with a vector subspace of \(G^{*}\) then the topology induced on F by \(\sigma(G^{*}, G)\) is \(\sigma(F, G)\) , and F is dense in \(G^{*}\) in the topology \(\sigma(G^{*}, G)\) (II, p.~43, cor. 4); from which the proposition follows.

Vector spaces that are complete for a weak topology are therefore the duals G* of arbitrary vector spaces G endowed with \(\sigma(\mathbf{G}^{*}, \mathbf{G})\) ; after II, p.~51, cor. 2, they are (topologically) isomorphic to products \(R^{I}\) of real lines. To simplify the language, we shall call them products of lines (for an intrinsic characterisation of these spaces see II, p.~85, exerc. 13 and II, p.~81, exerc. 1).

We note that on \(G^{*}\) , the \(\sigma(G^{*}, G)\) topology is minimal among the weak topologies that are Hausdorff; for, a weak topology that is coarser than \(\sigma(G^{*}, G)\) is necessarily of the form \(\sigma(G^{*}, H)\) where \(H \subset G\) (II, p.~43, cor. 3); but if \(H \neq G\) , then there exists a linear form \(x^{*} \in \mathbf{G}^{*}\) that is non null and is orthogonal to \(\mathbf{H}\) (A, II, § 7.3, prop. 8), therefore \(\sigma(\mathbf{G}^{*}, \mathbf{H})\) is not Hausdorff.

We deduce from this remark that, if F, G are two vector spaces, then a linear bijection \(u:G^{*}\to F^{*}\) , that is continuous for the topologies \(\sigma(G^{*},G)\) and \(\sigma(F^{*},F)\) , is necessarily bicontinuous.

PROPOSITION 10. --- Let G be a real vector space and F = G* its dual with the topology \(\sigma(\mathbf{G}^*, \mathbf{G})\) .

\begin{enumerate}
\def\labelenumi{(\roman{enumi})}
\item
  The mapping \(\mathbf{V} \mapsto \mathbf{V}^{\circ}\) is a bijection of the set of vector subspaces of \(\mathbf{G}\) on the set of closed vector subspaces of \(\mathbf{F}\) .
\item
  Every closed vector subspace of \(\mathbf{F}\) is a product of lines and has a topological complement.
\end{enumerate}

By the bipolar theorem (II, p.~45, cor. 2) \(\mathrm{V} \mapsto \mathrm{V}^{\circ}\) is a bijection of the set of vector subspaces V of G, closed for \(\sigma(\mathbf{G}, \mathbf{G}^{*})\) on the set of closed vector subspaces of F. But, by definition, every linear form on G is continuous for \(\sigma(\mathbf{G}, \mathbf{G}^{*})\) , therefore every vector subspace in G is closed, being defined by a system of equations \(\langle y, y_{\lambda}^{*} \rangle = 0\) (where \(y_{\lambda}^{*} \in \mathbf{G}^{*}\) ); this proves (i).

Now let W be a closed subspace of F; we have then \(W = V^{\circ}\) with \(V = W^{\circ}\) in G. Let \(V'\) be a complement of V in G. We know that \(F = G^{*}\) can be canonically identified with \(V^{*} \oplus V'^{*}\) , and \(V'^{*}\) identified with \(V^{\circ} = W\) (A, II, § 2.6, cor. to prop. 10); further (II, p.~50, prop. 8) the topology \(\sigma(G^{*}, G)\) can be identified with the product of the topologies \(\sigma(V^{*}, V)\) and \(\sigma(V'^{*}, V')\) ; this proves assertion (ii).

Though, for the topology \(\sigma(\mathbf{G},\mathbf{G}^{*})\) , every vector subspace of \(\mathbf{G}\) is closed, we note that if \(\mathbf{G}\) is of infinite dimension then the topology \(\sigma(\mathbf{G},\mathbf{G}^{*})\) is not the finest locally convex topology on \(\mathbf{G}\) , every neighbourhood of 0 for \(\sigma(\mathbf{G},\mathbf{G}^{*})\) containing a vector subspace of infinite dimension: it is however the finest of the weak topologies on \(\mathbf{G}\) (II, p.~43, cor. 3).

\subsection{Complete convex cones in weak spaces}

Lemma 1. --- Let E be a Hausdorff weak space and C a proper cone with vertex 0 in E, that is complete for the uniform structure induced by that of E. Every continuous linear form in E is then the difference between two continuous linear forms in E that are positive in C.

Let \(E'\) be the dual of E and F be the algebraic dual of \(E'\) , with the topology \(\sigma(F, E')\) . Let \(H = C^{\circ} - C^{\circ}\) be the vector subspace of \(E'\) formed by the differences of linear forms that are continuous in E and positive in C (II, p.~44, prop. 4). It is sufficient to show that the orthogonal to H in F is \(\{0\}\) (II, p.~41, Example 1). Then let \(a \in F\) be orthogonal to H; as a is orthogonal to \(C^{\circ}\) , it must belong to the bipolar of C in F. But E is identifiable as a subspace of F, and since C is complete, thus closed in F, we have \(a \in C\) (II, p.~44, th. 1). Similarly a is orthogonal to \(-C^{\circ}\) and therefore \(a \in -C\) . As C is proper, we have a = 0.

PROPOSITION 11. --- Let E be a Hausdorff weak space, and C be a proper convex cone with vertex 0 in E and which is complete in the uniform structure induced by that of E.

Then there exists a set I and a continuous linear mapping u of E in the product space \(R^{I}\) with the following properties :

\begin{enumerate}
\def\labelenumi{\alph{enumi})}
\item
  \(u\) is an isomorphism of \(\mathbf{C}\) on \(u(\mathbf{C})\) for the uniform structures induced respectively by those of \(\mathbf{E}\) and of \(\mathbf{R}^{\mathrm{I}}\) .
\item
  We have \(u(\mathbf{C}) \subset \mathbf{R}_{+}^{\mathrm{I}}\) .
\end{enumerate}

Further, if the uniform structure induced on C by that of E is metrisable, then we can take I = N.

Let \((f_{\iota})_{\iota \in I}\) be a family of continuous linear forms in E such that the finite sums of pseudometrics of the form \((x, y) \mapsto |f_{\iota}(x - y)|\) on \(\mathbf{C} \times \mathbf{C}\) define the uniform structure of C. (If the structure is metrisable we can take \(\mathbf{I} = \mathbf{N}\) .) By lemma 1 we can suppose further that each of the \(f_{\iota}\) is positive in C. Let \(u\) be the linear mapping \(x \mapsto (f_{\iota}(x))_{\iota \in I}\) of E in \(\mathbf{R}^{\mathrm{I}}\) . It is clear that \(u\) is continuous and that \(u(\mathbf{C}) \subset \mathbf{R}_{+}^{\mathrm{I}}\) . The restriction \(u|\mathbf{C}\) is a uniformly continuous mapping that is surjective from C on \(u(\mathbf{C})\) . Further if \(x, y\) in C are such that \(f_{\iota}(x) = f_{\iota}(y)\) for all \(\iota \in I\) , then \(x = y\) since the uniform structure of C is Hausdorff; therefore \(u|\mathbf{C}\) is bijective. Finally, if W is an entourage of the uniform structure of C, then there exists a finite set J of I and a number \(\varepsilon > 0\) such that the relations \(|f_{\iota}(x) - f_{\iota}(y)| \leqslant \varepsilon\) for \(\iota \in J\) imply \((x, y) \in W\) ; therefore \(u|\mathbf{C}\) is an isomorphism of C on \(u(\mathbf{C})\) for the uniform structures being considered.

COROLLARY 1. --- Let E be a Hausdorff weak space and C a proper convex cone of vertex 0 in E that is complete for the uniform structure induced by that of E. Then the mapping \((x, y) \mapsto x + y\) of \(C \times C\) in C is proper.

Because of prop. 11, we can suppose that \(\mathbf{E} = \mathbf{R}^{\mathbf{I}}\) and that \(\mathbf{C} = \mathbf{R}_{+}^{\mathbf{I}}\) (GT, I, § 10.1, cor. 1 and 4). But then the mapping \((x, y) \mapsto x + y\) of \(\mathbf{C} \times \mathbf{C}\) in \(\mathbf{C}\) is written as \(((\xi_{\mathfrak{i}}), (\eta_{\mathfrak{i}})) \mapsto (\xi_{\mathfrak{i}} + \eta_{\mathfrak{i}})\) , and we can restrict ourselves to proving that the continuous mapping \(f: (\xi, \eta) \mapsto \xi + \eta\) of \(\mathbf{R}_{+} \times \mathbf{R}_{+}\) in \(\mathbf{R}_{+}\) , is proper (GT, I, § 10.1, cor.3), Now, for all \(\zeta \in \mathbf{R}_{+}\) , we see that \(f(\zeta)\) is the set of pairs \((\xi, \zeta - \xi)\) such that \(0 \leqslant \xi \leqslant \zeta\) , therefore the inverse image by \(f\) of the interval \([0, \zeta]\) is the set of the \((\xi, \eta) \in \mathbf{R}_{+} \times \mathbf{R}_{+}\) such that \(\xi + \eta \leqslant \zeta\) , which is compact. The conclusion follows applying (GT, I, § 10.3, prop. 7).

COROLLARY 2. --- Let E be a Hausdorff weak space, and C a proper convex cone with vertex 0 in E, that is complete for the uniform structure induced by that of E.

\begin{enumerate}
\def\labelenumi{(\roman{enumi})}
\item
  For every point \(a\) of \(\mathbf{E}\) , the intersection \(\mathbf{C} \cap (a - \mathbf{C})\) is compact.
\item
  Let \(\mathbf{A},\mathbf{B}\) be two closed sets in C. Then \(\mathbf{A} + \mathbf{B}\) is a closed set in C.
\item
  The set of the \((x, y) \in \mathbf{C} \times \mathbf{C}\) such that \(x + y = a\) is compact from cor. 1 and from GT, I, § 10.2, th. 1, b). Now this set is also the set of the \((x, a - x)\) for \(x \in \mathbf{C} \cap (a - \mathbf{C})\) , which proves (i).
\item
  If A and B are closed in C, then \(A \times B\) is closed in \(C \times C\) , therefore \(A + B\) is closed in C after cor. 1 and GT, I, § 10.1, prop. 1.
\end{enumerate}

\section{EXTREMAL POINTS AND EXTREMAL GENERATORS}

\subsection{Extremal points of compact convex sets}

DEFINITION 1. --- Let A be a convex set in an affine space E. Then we say that a point \(x \in A\) is an extremal point of A if there does not exist an open segment that is contained in A and contains x.

In other words, the relations \(x=\lambda y+(1-\lambda)z\) , \(y\in A\) , \(z\in A\) , \(y\neq z\) and \(0\leqslant\lambda\leqslant1\) imply \(\lambda=0\) or \(\lambda=1\) (thus x=y or x=z). This implies that x cannot be the barycentre of a set of n points \(x_{i}\) of A carrying positive masses unless x is one of the \(x_{i}\) ; for this is just the definition when n=2; for arbitrary n argue by induction on n, as x is the barycentre of \(x_{1}\) and of the barycentre \(y_{1}\) of the \(x_{i}\) with \(2\leqslant i\leqslant n\) , therefore x is identical with \(x_{1}\) or \(y_{1}\) , and in the second case it is sufficient to apply the induction hypothesis.

To say that \(x\) is an extremal point of \(A\) also means that \(A - \{x\}\) is convex.

Examples. --- 1) In the space \(\mathbf{R}^n\) , all the points of the sphere \(\mathbf{S}_{n-1}\) are extremal points of the closed ball \(\mathbf{B}_n\) . For, if \(\sum_{i} y_i^2 \leqslant 1\) , \(\sum_{i} z_i^2 \leqslant 1\) and \(0 < \lambda < 1\) , the relation

\[
\lambda^ {2} \sum_ {i} y _ {i} ^ {2} + (1 - \lambda) ^ {2} \sum_ {i} z _ {i} ^ {2} + 2 \lambda (1 - \lambda) \sum_ {i} y _ {i} z _ {i} = 1 = (\lambda + (1 - \lambda)) ^ {2}
\]

is possible only if

\[
\sum_ {i} y _ {i} ^ {2} = \sum_ {i} z _ {i} ^ {2} = \sum_ {i} y _ {i} z _ {i} = 1.
\]

But this implies \(\sum_{i}(y_i - z_i)^2 = 0\) , thus \(y_i = z_i\) for all \(i\) , which proves our assertion.

\begin{enumerate}
\def\labelenumi{\arabic{enumi})}
\setcounter{enumi}{1}
\tightlist
\item
  In the normed space \(\mathcal{B}(\mathbf{N})\) of bounded sequences of real numbers (I, p.~4) the extremal points of the unit ball are the points \(x = (\xi_n)\) such that \(|\xi_n| = 1\) for all \(n\) . For, suppose that we had \(|\xi_n| \leqslant 1\) for all \(n\) and \(|\xi_p| < 1\) for one index \(p\) . We can then write
\end{enumerate}

\[
x = \frac {1 + \xi_ {p}}{2} y + \frac {1 - \xi_ {p}}{2} z
\]

where \(y\) (resp. \(z\) ) is the point all of whose coordinates are equal to the coordinate of \(x\) with the same index, except in the case of index \(p\) where the coordinate is equal to 1 (resp. - 1). This shows that \(x\) is not extremal, since we have \(\| y \| \leqslant 1\) and \(\| z \| \leqslant 1\) . Conversely, if \(|\xi_n| = 1\) for all \(n\) , then \(x\) is extremal, for the relation \(\xi_n = \lambda \eta_n + (1 - \lambda) \zeta_n\) with \(|\eta_n| \leqslant 1\) , \(|\zeta_n| \leqslant 1\) and \(0 < \lambda < 1\) implies \(\xi_n = \eta_n = \zeta_n\) .

\begin{enumerate}
\def\labelenumi{\arabic{enumi})}
\setcounter{enumi}{2}
\tightlist
\item
  Let \(u: \mathrm{E} \to \mathrm{E}'\) be an affine mapping of an affine space \(\mathrm{E}\) in an affine space \(\mathrm{E}'\) ; let \(\mathrm{C} \subset \mathrm{E}\) , \(\mathrm{C}' \subset \mathrm{E}'\) be two convex sets such that \(u(\mathrm{C}) \subset \mathrm{C}'\) . If \(x'\) is an extremal point of \(\mathrm{C}'\) and \(x\) is an extremal point of \(u^{-1}(x') \cap \mathrm{C}\) , then \(x\) is an extremal point of \(\mathrm{C}\) , as it follows from def. 1.
\end{enumerate}

PROPOSITION 1. --- Let B be the set of extremal points of A, a non-empty compact convex set in a Hausdorff locally convex space E, and let f be a convex function defined in A and upper semi-continuous. Then f attains its upper bound in A at one point (at least) of B.

Use F to denote the family of subsets X of A that are non-empty, closed, and such that every open segment that is contained in A and meets X necessarily lies in X. It has the following properties;

\begin{enumerate}
\def\labelenumi{(\roman{enumi})}
\item
  A belongs to \(\mathfrak{F}\) .
\item
  A point \(a \in \mathbf{A}\) is such that \(\{a\} \in \mathfrak{F}\) , if, and only if, \(a\) is an extremal point of \(\mathbf{A}\) .
\item
  Every non-empty intersection \(\mathbf{X}\) of a family \((\mathbf{X}_{\alpha})\) of sets of \(\mathfrak{F}\) also belongs to \(\mathfrak{F}\) .
\end{enumerate}

The properties (i), (ii) and (iii) follow immediately from the definitions.

\begin{enumerate}
\def\labelenumi{(\roman{enumi})}
\setcounter{enumi}{3}
\tightlist
\item
  Let \(X \in F\) , and let h be a function that is convex and upper semi-continuous in A; then the set Y of the points of X where the restriction \(h|X\) attains its upper bound in X is such that Y belongs to F.
\end{enumerate}

For, \(h|X\) being upper semi-continuous in \(X\) attains its upper bound \(\alpha\) over \(X\) in at least one point of \(X\) (GT, IV, § 6.2, th. 3); thus \(Y\) is non-empty, it is also closed (GT, IV, § 6.2, prop. 1). On the other hand let \(x, y\) be two distinct points of \(A\) and let \(z = \lambda x + (1 - \lambda)y\) be a point of \(Y\) such that \(0 < \lambda < 1\) ; as \(Y \subset X\) and \(X \in \mathfrak{F}\) , we have \(x \in X\) and \(y \in X\) ; on the other hand, as \(h\) is convex, we have

\[
h (z) \leqslant \lambda h (x) + (1 - \lambda) h (y)
\]

but as \(h(x) \leqslant \alpha\) , \(h(y) \leqslant \alpha\) and \(h(z) = \alpha\) , of necessity \(h(x) = h(y) = \alpha\) , that is to say \(x \in \mathbf{Y}\) and \(y \in \mathbf{Y}\) . Therefore \(\mathbf{Y} \in \mathfrak{F}\) .

With these properties established, let M be the set of \(x \in A\) where f attains its upper bound in A; by (iv), \(M \in F\) . On the other hand, by (iii) and the fact that the sets of F are closed subsets of the compact set A, it follows that F is inductive for the order relation \(\supset\) . By th. 2 of S, III, §2.4, M contains a subset N which is a minimal element of F. We shall show that N consists of a single point and this will complete the proof of the proposition. Since E is a Hausdorff locally convex space, it is sufficient to show that every continuous linear form u on E is constant in N (II, p.~38, cor. 1). Now it follows from (iv) that the set \(N'\) of the \(x \in N\) where \(u|N\) attains its upper bound in N is such that \(N'\) belongs to F; since N is minimal in F we necessarily have \(N' = N\) .

COROLLARY. --- Let A be a compact convex set in a Hausdorff locally convex space E. Then every closed support hyperplane H of A contains at least one extremal point of A.

For, if \(f(x) = \alpha\) is an equation of \(H\) and \(f(x) \leqslant \alpha\) in \(A\) , it is sufficient to apply prop. 1 to \(f\) .

THEOREM 1 (Krein-Milman). --- In a Hausdorff locally convex space E, every compact convex set A is the closed convex envelope of the set of its extremal points.

For, let C be the closed convex envelope of the set of extremal points of A; clearly \(C \subset A\) . To see that \(A \subset C\) , it is sufficient to prove that, if u is an affine linear function, continuous in E and if \(u(x) \geqslant 0\) in C then also \(u(x) \geqslant 0\) in A (II, p.~39, cor. 4); but this follows from prop. 1 applied to -u.

PROPOSITION 2. --- Let x be an extremal point of a compact convex set A in a Hausdorff locally convex space E. Then for every open neighbourhood V of x in E, there exists an open half-space F in E such that \(x \in F \cap A \subset V \cap A\) (in other words, the traces on A of the open half-spaces containing x, form a fundamental system of neighbourhoods of x in A).

For every open half-space D of E containing x, the set \(A \cap \overline{D}\) is a compact neighbourhood of x in A, and the intersection of all these neighbourhoods is precisely the point x (any two distinct points can be strictly separated by a closed hyperplane (II, p.~38, prop. 4). By prop. 1 of GT, I, § 9.2, it is sufficient to prove that the sets \(A \cap \overline{D}\) form a filter base. Now if we write \(L_{D} = A \cap (E - D)\) , the set \(L_{D}\) is convex, compact and contained in the convex set \(A - \{x\}\) ; if \(D_{1}, D_{2}\) are two open half-spaces of E containing x, the convex envelope B of \(L_{D_{1}} \cup L_{D_{2}}\) is therefore contained in \(A - \{x\}\) ; but B is a compact set (II, p.~14, prop. 15), therefore there exists a closed hyperplane H that separates x strictly from B (II, p.~38, prop. 4) and if the open half-space determined by H and containing x is D, then we have \(L_{D_{1}} \cup L_{D_{2}} \subset L_{D}\) , therefore \(A \cap \overline{D} \subset (A \cap \overline{D}_{1}) \cap (A \cap \overline{D}_{2})\) .

COROLLARY. --- In a Hausdorff locally convex space let K be a compact subset of a compact convex set A. Then the following conditions are equivalent.

\begin{enumerate}
\def\labelenumi{\alph{enumi})}
\item
  A is the closed convex envelope of K.
\item
  K meets every set that is the intersection of A with one of its support hyperplanes.
\item
  K contains the set of extremal points of A.
\end{enumerate}

\(a) \Rightarrow b)\) . Suppose that there exists a support hyperplane H of A whose equation is \(f(x) = \alpha\) , such that \((H \cap A) \cap K = \varnothing\) and suppose, for example, that \(f(x) \geqslant \alpha\) in A. As \(f(x) - \alpha > 0\) for all \(x \in K\) by hypothesis and as K is compact we have

\[
\beta = \inf _ {x \in \mathbf {K}} f (x) > \alpha ,
\]

and K is, therefore, contained in the closed half-space \(f(x) \geqslant \beta\) ; therefore the same is true of the closed convex envelope A of K and this is absurd.

\(b) \Rightarrow c)\) . Suppose that an extremal point \(x\) of \(\mathbf{A}\) does not belong to \(\mathbf{K}\) ; there is a neighbourhood \(\mathbf{V}\) of \(x\) in \(\mathbf{E}\) such that \(\mathbf{V} \cap \mathbf{A} \cap \mathbf{K} = \varnothing\) . But by prop. 2, we can suppose that \(\mathbf{V}\) is an open half-space defined by a hyperplane \(\mathbf{H}\) with the equation \(f(z) = \alpha\) . If for example \(f(x) > \alpha\) , then for all \(y \in \mathbf{K}\) , we have \(f(y) \leqslant \alpha\) , therefore \(\mathbf{K}\) does not meet the intersection of \(\mathbf{A}\) and the support hyperplane \(f(z) = \gamma > \alpha\) parallel to \(\mathbf{H}\) (II, p.~37, prop. 2); this is absurd.

\(c) \Rightarrow a)\) . This is an obvious consequence of the Krein-Milman theorem.

exerc. 11). The proof of theorem 1 (II, p.~56) shows that in any case A is the convex closed envelope of the set of extremal points of A which belong to a support hyperplane.

\subsection{Extremal generators of convex cones}

Let C be a convex cone with vertex 0 in a vector space E; clearly no other point of C than the vertex can be an extremal point; the vertex is an extremal point of C if and only if C is pointed and proper.

DEFINITION 2. --- Let C be a convex cone of vertex 0 in a vector space E. We say that a half-line D ⊂ C originating at 0 is an extremal generator of C, if every open segment contained in C, not containing 0 and meeting D is contained in D.

It comes to the same thing to say that for all \(x \in D\) such that \(x \neq 0\) , if \(y \neq 0\) , \(y' \neq 0\) are two points of C such that \(x = y + y'\) , then, it is necessarily the case that \(y \in D\) and \(y' \in D\) .

Remark 1. --- Let C be a pointed proper convex cone in E, and consider on E the order structure for which C is the set of elements \(\geqslant 0\) (II, p.~12, prop. 13); in order that an element of E, say x \textgreater{} 0, belongs to an extremal generator of C, it is necessary and sufficient that every element \(y \geqslant 0\) , that is bounded above by x, is of the form \(\lambda x\) with \(0 \leqslant \lambda \leqslant 1\) : in fact, to say that y is bounded above by x means that \(x = y + y'\) where \(y' \in C\) , whence the conclusion follows.

PROPOSITION 3. --- In a vector space E, let C be a convex cone with vertex 0, and let \(x_{0} \neq 0\) be a point of C, and D a half-line that is contained in C, originating from 0 and containing \(x_{0}\) . Let H be a hyperplane containing \(x_{0}\) and not passing through 0. Then D is an extremal generator of C if and only if \(x_{0}\) is an extremal point of H ∩ C.

The condition is clearly necessary. Conversely, suppose that it is satisfied; suppose that there is a line \(D'\) not containing D, passing through \(x_{0}\) and such that \(D' \cap C\) contains an open segment to which \(x_{0}\) belongs. Let \(y \neq 0\) be a direction vector of \(D'\) ; the hypotheses imply that the point \((1 + \lambda)x_{0} + \mu y\) belongs to C for \(|\lambda|\) and \(|\mu|\) sufficiently small. But then, in the plane P determined by D and \(D'\) and carrying the canonical topology, \(x_{0}\) is an interior point of \(P \cap C\) , and it follows that the line \(P \cap H\) contains an open segment contained in \(H \cap C\) and to which \(x_{0}\) belongs. This contradicts the hypothesis.

DEFINITION 3. --- Let C be a convex set in a Hausdorff topological vector space E. A compact convex non-empty set A of C is called a cap of C if the complement C−A of A in C is convex.

Let C be a pointed convex cone with vertex 0 in E and let A be a cap of C. Write B = C - A. For every closed half-line L ⊂ C originating at 0, the sets L ∩ A and L ∩ B are convex sets that are complements in L, whose union is L, and such that L ∩ A is compact. As L ∩ A is non-empty for at least one half-line L, we see that 0 ∈ A, thus L ∩ A is a closed segment with an end point at 0. If A exists then C is proper.

PROPOSITION 4. --- Let C be a pointed convex cone with vertex 0 in E, a Hausdorff locally convex space.

\begin{enumerate}
\def\labelenumi{\alph{enumi})}
\tightlist
\item
  Let A be a cap of C. Let p be the restriction to C of the gauge of A (II, p.~20). The set of the \(x \in C\) such that \(p(x) \leqslant 1\) is the set A. The function p is lower semicontinuous and has the following properties :
\end{enumerate}

\begin{enumerate}
\def\labelenumi{(\roman{enumi})}
\item
  For any \(x, y\) in \(\mathbf{C}\) , we have \(p(x + y) = p(x) + p(y)\) .
\item
  For any \(x \in \mathbf{C}\) and \(\lambda \in \mathbf{R}_{+}^{*}\) , we have \(p(\lambda x) = \lambda p(x)\) .
\item
  If \(x \in C\) , the relation \(p(x) = 0\) is equivalent to x = 0.
\end{enumerate}

\begin{enumerate}
\def\labelenumi{\alph{enumi})}
\setcounter{enumi}{1}
\tightlist
\item
  Conversely, let \(p\) be a function defined in \(\mathbf{C}\) with values in \([0, +\infty]\) , satisfying the conditions (i), (ii) of a). Let \(\mathbf{A}\) be the set of the \(x \in \mathbf{C}\) such that \(p(x) \leqslant 1\) . Then \(\mathbf{A}\) and \(\mathbf{C} - \mathbf{A}\) are convex. A is a cap, if and only if \(\mathbf{A}\) is compact and non-empty.
\end{enumerate}

The statement \(b)\) is obvious. The properties stated in \(a)\) are consequences of the remarks preceding prop. 4 and of the prop. 22 of II, p.~20 and prop. 23 of II, p.~20 with the exception of

\[
p (x + y) \geqslant p (x) + p (y).
\]

It is sufficient to prove this last when \(x \neq 0\) and \(y \neq 0\) ; we have therefore \(p(x) > 0\) , \(p(y) > 0\) . Let \(\mu, \lambda\) be two numbers \textgreater{} 0 such that \(\lambda < p(x)\) , \(\mu < p(y)\) , and denote the complement of A in C by B. We have \(x \in \lambda B\) , \(y \in \mu B\) , therefore \(x + y \in \lambda B + \mu B\) ; by the convexity of B, we have \(\lambda B + \mu B \subset (\lambda + \mu)B\) , whence \(p(x + y) > \lambda + \mu\) , which implies the inequality stated above.

COROLLARY 1. --- Let C be a pointed convex cone of vertex 0 in E, a Hausdorff locally convex space and let p be the gauge of A, a cap of C. The extremal points of A are then the point 0, and the points x on the extremal generators of C such that \(p(x) = 1\) .

It is clear that 0 is an extremal point of A. Let x be a point on L an extremal generator of C and such that \(p(x) = 1\) . Let y, z be two points of A such that \(x = \frac{1}{2}(y + z)\) . As L is extremal, we have \(y = \lambda x\) and \(z = \mu x\) , where \(\lambda\) and \(\mu\) are numbers \(\geqslant 0\) such that \(\frac{1}{2}(\lambda + \mu) = 1\) , \(\lambda = \lambda p(x) = p(y) \leqslant 1\) and \(\mu = \mu p(x) = p(z) \leqslant 1\) , from which \(\lambda = \mu = 1\) and hence y = z = x; so, x is an extremal point of A. Conversely, let \(x \neq 0\) be an extremal point of A. Obviously \(p(x) = 1\) . Let y, \(y'\) be two points of C such that \(x = y + y'\) , and we shall show that y, \(y'\) are proportional to x. Without loss of generality we can suppose that the numbers \(\lambda = p(y)\) and \(\lambda' = p(y')\) are finite and \textgreater0. Then \(\lambda^{-1}y \in A\) , \(\lambda'^{-1}y' \in A\) , \(\lambda + \lambda' = 1\) by prop. 4, (i) and the equality \(x = \lambda(\lambda^{-1}y) + \lambda'(\lambda'^{-1}y')\) implies, by hypothesis that

\[
x = \lambda^ {- 1} y = \lambda^ {\prime - 1} y ^ {\prime}.
\]

COROLLARY 2. --- Every point of C that belongs to a cap of C, also belongs to the convex closed envelope of the union of the extremal generators of C.

This follows immediately from cor. 1 and the Krein-Milman theorem (II, p.~55, th. 1).

* Example. --- Let X be a locally compact space that is \(\sigma\) -compact. Let C be a closed convex cone of vertex 0 in \(\mathcal{M}_{+}(X)\) with the vague topology. We shall show that C is the union of its caps. Let \((X_{n})\) be an increasing sequence of open, relatively compact sets of X whose union is X. Let \(\mu\) be an element \(\neq 0\) of C. There exist \(\alpha_{n} > 0\) such that \(\sum \alpha_{n}\mu(X_{n}) = 1\) .

For every measure \(v \in C\) , put \(p(v) = \sum_{n} \alpha_{n} v(X_{n}) \in [0, +\infty]\) . The function p on C satisfies conditions (i) and (ii) of prop. 4. It is lower semi-continuous for the vague topology (INT, IV, 2nd ed., § 1, No.~1, prop. 4). The set A of the \(\gamma \in C\) such that \(p(\gamma) \leqslant 1\) is therefore closed and non-empty. On the other hand, every compact set of X is contained in one of the \(X_{n}\) , thus A being vaguely bounded is also vaguely compact (INT, III, 2nd ed., § 1, No.~9, prop. 15). The set A is therefore a cap of C containing \(\mu\) .

PROPOSITION 5. --- Let C be a proper convex cone with vertex 0 in E, a Hausdorff weak space; suppose that C is complete for the uniform structure induced by that of E, and that there is an enumerable fundamental system of neighbourhoods of 0 in C. Then C is the union of its caps and is the closed convex envelope of the union of its extremal generators.

The second statement follows from the first and from cor. 2 above. Using prop. 11 of II, p.~52 reduces the proposition to the case when \(\mathbf{E} = \mathbf{R}^{\mathrm{I}}\) and \(\mathbf{C} \subset \mathbf{R}_{+}^{\mathrm{I}}\) . For all \(\alpha \in \mathbf{I}\) , denote the projection \(\operatorname{pr}_{\alpha}\) in \(\mathbf{E}\) by \(f_{\alpha}\) ; then \(f_{\alpha}\) is a continuous linear form. On the other hand let \((\mathrm{V}_n)_{n \in \mathbb{N}}\) be an enumerable fundamental system of neighbourhoods of 0 in \(\mathbb{C}\) . By the definition of the topology of \(\mathbf{E}\) , for each \(n \in \mathbb{N}\) , there exists a finite subset \(\mathrm{J}_n\) of \(\mathbf{I}\) and a number \(\varepsilon_n > 0\) such that \(\mathrm{V}_n\) contains the set \(\mathrm{W}_n\) of the \(x \in \mathbb{C}\) such that \(f_{\alpha}(x) \leqslant \varepsilon_n\) for all \(\alpha \in \mathrm{J}_n\) ; put \(\mathrm{J} = \bigcup_{n \in \mathbb{N}} \mathrm{J}_n\) . Let \(y \neq 0\) be a point of \(\mathbb{C}\) , and \(p\) be the function \(\sum_{\alpha \in \mathbf{J}} \lambda_{\alpha}(f_{\alpha}|\mathbb{C})\) where the \(\lambda_{\alpha} > 0\) are chosen so that \(p(y) = 1\) ; this is possible, since if \(f_{\alpha}(y) = 0\) for all \(\alpha \in \mathbf{J}\) , then \(y \in \mathrm{V}_n\) for all \(n\) , which implies that \(y = 0\) , and this is contrary to hypothesis. Now we remark that for all \(\alpha \in \mathbf{I}\) , the function \(f_{\alpha}|\mathbb{C}\) is continuous at the point 0, therefore there is an \(n \in \mathbb{N}\) , such that \(f_{\alpha}\) is bounded in a \(\mathrm{W}_n\) , therefore bounded above in \(\mathbb{C}\) by a linear combination of a finite number of functions \(f_{\beta}|\mathbb{C}\) , where \(\beta \in \mathbf{J}\) . It follows that if \(\mathbf{A}\) in the set of \(x \in \mathbb{C}\) such that \(p(x) \leqslant 1\) , then \(f_{\alpha}\) is bounded in \(\mathbf{A}\) for all \(\alpha \in \mathbf{I}\) . As \(p\) is lower semi-continuous in \(\mathbb{C}\) , it follows that \(\mathbf{A}\) is closed and non-empty in \(\mathbb{C}\) and therefore is compact. Since it is clear that \(p\) verifies the conditions (i) and (ii) of prop. 4 of II, p.~58, we see that \(\mathbf{A}\) is a cap in \(\mathbb{C}\) and contains \(y\) .

Remark 2. --- There exist proper convex cones that are weakly complete and which have no extremal generator (II, p.~92, exerc. 31).

\subsection{Convex cones with compact sole}

PROPOSITION 6. --- Let E be a Hausdorff locally convex space and K a convex compact set in E which does not contain 0. Then the smallest pointed cone C of vertex 0 which contains K is a proper convex cone in E and is a locally compact and complete subspace of E; also, there exists a closed hyperplane H in E that does not contain 0 and is such that H meets all the half-lines originating at 0 contained in C and such that H ∩ C is compact. Further, if D is the half-space containing 0 determined by H, a closed hyperplane with these properties, then C ∩ D is a cap of C and C is the union of the λ(C ∩ D) for λ \textgreater{} 0.

By prop. 4 of II, p.~38, there exists a closed hyperplane H which separates 0 strictly from K. Now, the convex envelope A of the union of \(\{0\}\) and of K is compact (II, p.~14, prop. 15) and is the union of the \(\lambda K\) with \(0 \leqslant \lambda \leqslant 1\) . As 0 and K are strictly on opposite sides of H, for every \(x \in K\) there exists \(\lambda\) such that \(0 < \lambda < 1\) and \(\lambda x \in H\) . As C is the union of the \(\lambda A\) for \(\lambda \geqslant 1\) , we see that H meets every half-line originating at 0 contained in C and that \(H \cap A = H \cap C\) is compact. Further, C is also the union of the \(\lambda(H \cap C)\) for \(\lambda \geqslant 0\) ; let \(C_n\) be the union of the \(\lambda(H \cap C)\) for \(0 \leqslant \lambda \leqslant n\) . Clearly \(C_n\) is the convex envelope of the union of \(\{0\}\) and of \(n(H \cap C)\) , therefore it is compact. Also, for all \(x \in E\) , there is a closed neighbourhood V of x in E and an integer n such that \(V \cap C \subset C_n\) ; in fact, if H is defined by the equation \(f(z) = \alpha\) , where \(\alpha > 0\) , it is sufficient to take for V the closed half-space determined by nH and containing 0, where n is so large that \(n\alpha > f(x)\) . This shows that C is locally compact (taking \(x \in C\) ), and that it is closed in E. We can also consider K as a subset of the completion \(\hat{E}\) , therefore C is also closed in \(\hat{E}\) and therefore complete.

Given a cone C and a closed hyperplane H in a Hausdorff topological vector space E, such that H does not contain the vertex s of C and C is the smallest cone with vertex s containing \(H \cap C\) , then we call the intersection \(H \cap C\) a « sole » of the cone C. Prop. 6 shows that in a Hausdorff locally convex space E, the smallest cone of vertex 0, containing a compact convex set K to which 0 does not belong, is a cone of compact sole, and that every convex cone having a compact sole S, is locally compact and complete.

Examples. --- 1) Every proper closed convex cone in E, a vector space of finite dimension, has a compact sole. In fact, by II, p.~52, prop. 11 we need only consider the case where \(E = R^{n}\) and \(C = R_{+}^{n}\) . If \((e_{i})_{1 \leqslant i \leqslant n}\) is the canonical basis of \(R^{n}\) , it is clear that the compact convex set which is the convex envelope of the \(e_{i}\) ( \(1 \leqslant i \leqslant n\) ) is a compact sole for \(R_{+}^{n}\) .

* 2) If X is a compact space, then the cone \(\mathcal{M}_{+}(\mathrm{X})\) of positive measures on X, with the vague topology, is a cone with a compact sole (INT, III, 2nd ed., § 1, No.~9, cor. 3 of prop. 15). *

\section{COMPLEX LOCALLY CONVEX SPACES}

\subsection{\texorpdfstring{Topological vector spaces over \(\mathbf{C}\)}{Topological vector spaces over the complex field}}

Let E be a topological vector space over C the field of complex numbers; the topology of E is also compatible with the structure of the vector space over R, obtained by restricting the field of scalars to R. We denote by \(E_{0}\) the topological vector space on R which underlies E (I, p.~2). Note that, in \(E_{0}\) , the mapping \(x \mapsto ix\) (which is not a homothety) is an automorphism u of the topological vector space structure of \(E_{0}\) such that \(u^{2}(x) = -x\) .

Conversely, let \(\mathbf{F}\) be a topological vector space over \(\mathbf{R}\) , and suppose that there exists an automorphism \(u\) of \(\mathbf{F}\) such that \(u^2 = -1_{\mathbf{F}}(1_{\mathbf{F}}\) is the identity automorphism of \(\mathbf{F}\) . We know (A, IX, § 3, No.~2) that it is then possible to define a vector space structure on \(\mathbf{F}\) relative to \(\mathbf{C}\) writing \(\lambda x = \alpha x + \beta u(x)\) for all \(\lambda = \alpha + i\beta \in \mathbf{C}\) and all \(x \in \mathbf{F}\) . Further since the mapping \((\alpha, \beta, x) \mapsto \alpha x + \beta u(x)\) of \(\mathbf{R}^2 \times \mathbf{F}\) in \(\mathbf{F}\) is continuous the topology of \(\mathbf{F}\) is compatible with the vector space structure relative to \(\mathbf{C}\) defined above; if \(\mathbf{E}\) denotes the topological vector space on \(\mathbf{C}\) defined in this manner, then \(\mathbf{F}\) is the topological vector space on \(\mathbf{R}\) which underlies \(\mathbf{E}\) .

Remark. --- Given a topological vector space F over R, it is not always the case that there exists an automorphism u of F whose square is \(-1_{F}\) ; for example, it is not possible to define vector space structure relative to C on a vector space over R of finite odd dimension.

Let E be a topological vector space on C, and \(E_{0}\) the topological vector space on R which underlies E. Every linear variety M in E is also a linear variety in \(E_{0}\) , but the converse is false. To avoid confusion we say that a linear variety for a vector space structure relative to C (resp. relative to R) is a complex (resp. real) linear variety. A complex linear variety of finite dimension n (resp. of finite codimension n) is a real linear variety of dimension 2n (resp. of codimension 2n). In order that a real vector subspace M of E should also be a complex vector subspace, it is necessary and sufficient that \(iM \subset M\) .

Recall that, if E and F are two topological vector spaces on C, then a mapping of E in F is called C-linear (resp. R-linear) if it is a linear mapping for the vector space structures of E and of F relative to C (resp. R); every C-linear mapping is evidently R-linear but the converse is false. We say that a C-linear form on E is a complex linear form and that an R-linear form on E (i.e.~a linear form on \(E_{0}\) ) is a real linear form. If f is a complex linear form on E, it is clear that the real part g = Rf and the imaginary part h = If of f are real linear forms; further, the relation \(f(ix) = if(x)\) implies the identity \(h(x) = -g(ix)\) ; in other words we have

\[
f (x) = (\mathscr {R} f) (x) - i (\mathscr {R} f) (i x).\tag{1}
\]

Conversely, if g is a real linear form on E, then \(f(x) = g(x) - ig(ix)\) is the unique complex linear form on E such that Rf = g; and f is continuous if, and only if, g is continuous.

Now let H be a complex hyperplane in E, with the equation \(f(x) = \alpha + i\beta\) , where f is a complex linear form on E; putting \(g = Rf\) , we see that H is the intersection of two real hyperplanes \(H_{1}\) , \(H_{2}\) with equations respectively \(g(x) = \alpha\) and \(g(ix) = -\beta\) ; if H is closed, so also are \(H_{1}\) and \(H_{2}\) (I, p.~13, th. 1). Conversely let \(H_{0}\) be a homogeneous real hyperplane, with equation \(g(x) = 0\) (where g is a real linear form on E); then H, the intersection of \(H_{0}\) and \(iH_{0}\) , is a homogeneous complex hyperplane, and if f is the complex linear form such that Rf = g, then \(f(x) = 0\) is the equation of H; if \(H_{0}\) is closed then H also is closed.

Let G be a topological vector space over R and let \(\mathbf{G}_{(\mathbf{C})}\) be the vector space on C obtained from G by extending the field of scalars to C (A, II, § 5.1). Identify G as a subset of \(\mathbf{G}_{(\mathbf{C})}\) by the mapping \(x \mapsto 1 \otimes x\) . The R-linear mapping \((x, y) \mapsto x + i \cdot y\) is then a bijection of \(G \times G\) on \(\mathbf{G}_{(\mathbf{C})}\) , by means of which we transfer the product topology of \(G \times G\) to \(\mathbf{G}_{(\mathbf{C})}\) . Then \(\mathbf{G}_{(\mathbf{C})}\) with this topology is a topological vector space on C. We say that \(\mathbf{G}_{(\mathbf{C})}\) is the complexified topological vector space of G.

\subsection{Complex locally convex spaces}

To say that a subset A of a complex vector space E is balanced means that, for all \(x \in A\) , we have \(\rho x \in A\) for \(0 \leqslant \rho \leqslant 1\) and \(e^{i\vartheta}x \in A\) for all real \(\vartheta\).

We say that a set A of E is convex if it is convex in the real space \(E_{0}\) which underlies E. In order that a convex set \(A \neq \varnothing\) of E be balanced, it is sufficient that \(e^{i\vartheta}A \subset A\) for all real \(\vartheta\) ; for this implies firstly that -A = A; as A is convex, we see that 0 belongs to A and thus \(\rho A \subset A\) for \(0 \leqslant \rho \leqslant 1\) .

Let E be a complex topological vector space. The smallest balanced convex (resp. closed balanced convex) set containing a set A of E is called the balanced convex envelope (resp. balanced closed convex envelope) of A; the balanced closed convex envelope of A is the closure of the balanced convex envelope of A. This last is the convex envelope of the union of the sets \(e^{i\theta}A\) ; we can therefore define it as the set of linear sums \(\sum_{i}\lambda_{i}x_{i}\) , when \((x_{i})\) is any finite family of points of A, and \((\lambda_{i})\) a family of complex numbers such that \(\sum_{i}|\lambda_{i}| \leqslant 1\) . If A is precompact so also is its balanced envelope (I, p.~6, prop. 3).

We say that a complex topological vector space E is locally convex if the real underlying topological vector space \(E_{0}\) is locally convex, that is to say if every neighbourhood of 0 in E contains a convex neighbourhood of 0; a topology T on E is locally convex if it is compatible with the vector space structure of E (relative to C) and if E, with topology T, is locally convex. As in this case every closed convex neighbourhood V of 0 contains a balanced neighbourhood W of (I, p.~7, prop. 4), we see that V also contains U, the balanced closed convex envelope of W; in other words the balanced, closed, convex neighbourhoods of 0 form a fundamental system of neighbourhoods of 0 in E, invariant under every homothety of ratio \(\neq 0\) .

Conversely, let E be a complex vector space and let S be a filter base on E formed by absorbent, balanced convex sets. We know then (II, p.~23, prop. 1) that the set B, of the transforms of the sets of S by homotheties of ratio \textgreater{} 0, is a fundamental system of neighbourhoods of 0 for a locally convex topology T on the real vector space \(E_{0}\) underlying E. Further, as the sets of B are balanced, they are invariant under every homothety \(x \mapsto e^{i\vartheta}x\) , which shows that T is compatible with the vector space structure of E (over C) (I, p.~7, prop. 4).

Every locally convex topology on a complex vector space E can be defined by a set of semi-norms, for the gauge of an open balanced convex neighbourhood of 0 is a semi-norm on E.

The ideas and results for real locally convex spaces detailed in II, p.~25 to 36, extend to complex locally convex spaces with no modification other than the replacement of symmetric convex sets by balanced convex sets.

A complex locally convex space is a Fréchet space if it is metrisable and complete.

\subsection{The Hahn-Banach theorem and its applications}

THEOREM 1 (Hahn-Banach). --- Let V be a vector subspace of E, a complex vector space, and let f be a (complex) linear form on V and p a semi-norm on E such that \(|f(y)| \leqslant p(y)\) for all \(y \in V\) . Then there exists a linear form \(f_{1}\) on E extending f and such that \(|f_{1}(x)| \leqslant p(x)\) for all \(x \in E\) .

For \(g = \mathcal{R}f\) is a real linear form defined in \(V\) and satisfying \(|g(y)| \leqslant p(y)\) at every point of \(V\) ; therefore there exists a real linear form \(g_1\) in \(E\) extending \(g\) and such that \(|g_1(x)| \leqslant p(x)\) for all \(x \in E\) (II, p.~23, cor. 1). Let \(f_1(x) = g_1(x) - ig_1(ix)\) be the complex linear form on \(E\) of which \(g_1\) is the real part (II, p.~61). For all real \(\vartheta\)

\[
\left| \mathcal {R} (e ^ {i \vartheta} f _ {1} (x)) \right| = \left| \mathcal {R} (f _ {1} (e ^ {i \vartheta} x)) \right| = \left| g _ {1} (e ^ {i \vartheta} x) \right| \leqslant p (e ^ {i \vartheta} x) = p (x)
\]

since \(p\) is a semi-norm on the complex space \(\mathbf{E}\) ; this implies the relation \(|f_1(x)| \leqslant p(x)\) , and the theorem is proved.

COROLLARY 1. --- Let \(x_0\) be a point of a complex topological vector space \(E\) and \(p\) be a continuous semi-norm in \(E\) ; then there exists a continuous (complex) linear form \(f\) defined in \(E\) , such that \(f(x_0) = p(x_0)\) and \(|f(x)| \leqslant p(x)\) for all \(x \in E\) .

COROLLARY 2. --- Let V be a vector subspace of a complex locally convex space E and f be a (complex) linear form defined and continuous in V; then there exists a continuous linear form \(f_{1}\) defined in E and extending f.~If E is normed there exists such a form \(f_{1}\) that also satisfies \(\|f_{1}\|=\|f\|\) .

COROLLARY 3. --- Let M be a finite dimensional vector subspace of a Hausdorff complex locally convex space E. Then there exists a closed vector subspace N of E that is a topological complement of M in E.

The proofs using theorem 1, p.~24 are the same as those of II, p.~23, cor. 2 and cor. 3, p.~24, prop. 2 and p.~25, cor. 2.

PROPOSITION 1. --- Let A be an open non-empty convex set in a complex topological vector space E and M be a non-empty (complex) linear variety that does not meet A. Then there exists a closed complex hyperplane H that contains M and does not meet A.

We can suppose that \(0 \in M\) . Then there exists a closed real hyperplane \(H_{0}\) containing M and not meeting A (II, p.~36; th. 1). As M = iM, the closed complex hyperplane \(\mathbf{H} = \mathbf{H}_{0} \cap (i\mathbf{H}_{0})\) has the properties required.

COROLLARY. --- In a complex locally convex space E, every closed complex linear variety M is the intersection of the closed complex hyperplanes which contain it.

In fact, for all \(x \notin M\) , there exists a convex open neighbourhood V of x that does not meet M, and thus there exists a closed complex hyperplane H containing M and not meeting V; a fortiori H does not contain x.

PROPOSITION 2. --- Let A be a non-empty balanced open convex set of a complex topological vector space E, and B be a non-empty convex set that does not meet A. Then there exists a continuous complex linear form f on E and a number \(\alpha > 0\) such that \(|f(x)| < \alpha\) in A and \(|f(y)| \geqslant \alpha\) in B.

For, there exists a continuous real linear form \(g\) on \(\mathbf{E}\) and a real number \(\alpha\) such that \(g(x) < \alpha\) in \(\mathbf{A}\) and \(g(y) \geqslant \alpha\) in \(\mathbf{B}\) (II, p.~37, prop. 1). As \(0 \in \mathbf{A}\) , we have \(\alpha > 0\) . We show that the continuous complex linear form \(f(x) = g(x) - ig(ix)\) and the number \(\alpha\) have the properties required. For, since \(\mathcal{R}f = g\) , we have \(|f(y)| \geqslant \alpha\) in \(\mathbf{B}\) . On the other hand, for all \(x \in \mathbf{A}\) and all real \(\vartheta\) , the point \(e^{i\vartheta}x\) belongs to \(\mathbf{A}\) , since \(\mathbf{A}\) is balanced, and we have \(f(x) = e^{-i\vartheta}f(e^{i\vartheta}x)\) ; then there exists a number \(\vartheta\) such that \(|f(x)| = \mathcal{R}(e^{i\vartheta}f(x)) = g(e^{i\vartheta}x) < \alpha\) , and the proposition follows.

PROPOSITION 3. --- Let A be a balanced, closed, convex set in a complex locally convex space E and let K be a non-empty compact convex set in E that does not meet A. Then there exists a continuous complex linear form f on E and a number \(\alpha > 0\) such that \(|f(x)| < \alpha\) in A and \(|f(y)| > \alpha\) in K.

The proposition follows from II, p.~38, prop. 4 as prop. 2 follows from II, p.~37, prop. 1.

\subsection{Weak topologies on complex vector spaces}

The definition and results of II, § 6, Nos. 1 and 2 apply without change to complex vector spaces. If F and G are two complex vector spaces in duality by a bilinear form B, then the underlying spaces \(F_{0}\) and \(G_{0}\) are in duality by RB, and it follows from II, p.~61, formula (1) that the weak topologies \(\sigma(\mathrm{F},\mathrm{G})\) and \(\sigma(\mathrm{F}_{0},\mathrm{G}_{0})\) are identical.

DEFINITION 1. --- Let F and G be two complex vector spaces in duality. For any subset M of F, the polar of M in G, denoted by \(M^{\circ}\) , is the set of \(y \in G\) such that \(\mathcal{R}(\langle x, y \rangle) \geqslant -1\) for all \(x \in M\) .

If \(\mathbf{M}^{\circ}\) is the polar of \(\mathbf{M} \subset \mathbf{F}\) in \(\mathbf{G}\) then \((\lambda \mathbf{M})^{\circ} = \lambda^{-1}\mathbf{M}^{\circ}\) for all \(\lambda \in \mathbf{C}^{*}\) .

If M is a (complex) vector subspace of E, then \(M^{\circ}\) is a closed vector subspace (for \(\sigma(\mathbf{G},\mathbf{F})\) ), since the relation \(\mathcal{R}(\lambda\langle x,y\rangle)\geqslant-1\) for every scalar \(\lambda\in C\) implies \(\langle x,y\rangle=0\) ; again we say that \(M^{\circ}\) is the subspace of G orthogonal to M.

If M is a balanced set in F, then \(M^{\circ}\) is a balanced set in G; in this case \(M^{\circ}\) is the set of \(y \in G\) such that \(\left|\langle x, y \rangle\right| \leqslant 1\) for all \(x \in M\) ; for this relation is equivalent to \(\mathcal{R}(\langle \zeta x, y \rangle) \leqslant 1\) for all \(x \in M\) and all \(\zeta \in C\) such that \(|\zeta| = 1\) .

The results of II, p.~41 to 51 are also valid without restriction for complex vector spaces.

\exercisesheading
